\subsection{Questão 85}

\subsubsection{Enunciado}

Um feixe de luz não-polarizada atravessa dois filtros polarizadores.

Qual deve ser o ângulo entre as direções de polarização dos filtros para que a intensidade da luz que atravessa os dois filtros seja um terço da intensidade da luz incidente?


\subsubsection{Resolução}

Ao atravessar o primeiro polarizador, a luz não-polarizada tem sua intensidade
reduzida à metade:

\[
I_1=\frac{I_0}{2}.
\]

Ao atravessar o segundo polarizador, vale a lei de Malus:

\[
I_2=I_1\cos^2\theta.
\]

Queremos

\[
I_2=\frac{I_0}{3}.
\]

Assim,

\[
\frac{I_0}{3}
=
\frac{I_0}{2}\cos^2\theta.
\]

Cancelando \(I_0\),

\[
\cos^2\theta=\frac{2}{3}.
\]

Logo,

\[
\theta=\cos^{-1}\left(\sqrt{\frac{2}{3}}\right)
=35{,}3^\circ.
\]

Portanto, o ângulo entre as direções de polarização dos filtros deve ser

\[
\boxed{35{,}3^\circ}.
\]

\vspace{1cm}
